% Manual de TurbuLENS 1.4.1 — compilar con XeLaTeX (o LuaLaTeX), dos pasadas.
% Overleaf: Menú → Compilador → XeLaTeX. Las imágenes van en la carpeta img/.
\documentclass[12pt,a4paper]{article}

\usepackage{fontspec}
\usepackage{polyglossia}
\setdefaultlanguage{spanish}
\IfFontExistsTF{Inter}{\setmainfont{Inter}[Scale=0.92]}{\IfFontExistsTF{Roboto}{\setmainfont{Roboto}}{}}
\IfFontExistsTF{DejaVu Sans Mono}{\setmonofont{DejaVu Sans Mono}[Scale=0.85]}{}
\usepackage{amsmath,amssymb}
\usepackage[margin=2.3cm,headheight=15pt]{geometry}
\usepackage{graphicx,xcolor,booktabs,tabularx,array,float,enumitem,multicol,caption,fancyhdr,titlesec}
\usepackage[most]{tcolorbox}
\usepackage{ragged2e}
\usepackage[hidelinks]{hyperref}

\definecolor{acento}{HTML}{3E7C72}
\definecolor{acento2}{HTML}{B4652A}
\definecolor{naranja}{HTML}{F28C28}
\definecolor{suave}{HTML}{E3EEEC}
\definecolor{gris}{HTML}{55615B}
\definecolor{aviso}{HTML}{9A6B12}
\hypersetup{colorlinks=true,linkcolor=acento,urlcolor=acento,citecolor=acento}

\setlength{\parindent}{0pt}
\setlength{\parskip}{0.6em}
\RaggedRight
\setlist{itemsep=0.15em,topsep=0.3em}
\newcolumntype{Y}{>{\RaggedRight\arraybackslash}X}
\captionsetup{font=small,labelfont={bf,color=acento}}
\titleformat{\section}{\Large\bfseries\color{acento}}{\thesection}{0.6em}{}
\titleformat{\subsection}{\large\bfseries}{\thesubsection}{0.6em}{}
\titleformat{\subsubsection}{\normalsize\bfseries\color{gris}}{\thesubsubsection}{0.6em}{}
\pagestyle{fancy}\fancyhf{}
\fancyhead[L]{\small\color{gris}Manual de TurbuLENS 1.4.1}
\fancyhead[R]{\small\color{gris}\leftmark}
\fancyfoot[C]{\small\thepage}
\renewcommand{\headrulewidth}{0.4pt}
\renewcommand{\sectionmark}[1]{\markboth{#1}{}}

\newcommand{\codigo}[1]{\texttt{#1}}
\newcommand{\boton}[1]{\textbf{#1}}
\newtcolorbox{nota}[1]{colback=suave,colframe=acento,boxrule=0.6pt,arc=2mm,left=3mm,right=3mm,top=2mm,bottom=2mm,
  fonttitle=\bfseries,coltitle=acento,colbacktitle=suave,title={#1},attach title to upper={\ \ }}
\newtcolorbox{cuidado}[1]{colback=orange!6,colframe=acento2,boxrule=0.6pt,arc=2mm,left=3mm,right=3mm,top=2mm,bottom=2mm,
  fonttitle=\bfseries,coltitle=acento2,title={#1},attach title to upper={\ \ }}
\newcommand{\captura}[2]{\begin{figure}[H]\centering\includegraphics[width=\textwidth]{img/#1}\caption{#2}\end{figure}}

\begin{document}

% ===================================================================================== portada
\begin{titlepage}
\centering
\vspace*{2.2cm}
\includegraphics[width=3.2cm]{img/logo.png}\par\vspace{1cm}
{\fontsize{40}{44}\selectfont\color{acento} TurbuLENS}\par\vspace{0.4cm}
{\Large Manual de usuario y fundamentos}\par\vspace{0.3cm}
{\large\color{gris} Versión de la aplicación 1.4.1}\par\vspace{1.4cm}
\begin{minipage}{0.82\textwidth}\centering\color{gris}
Medición del \emph{seeing} con un DIMM · enseñanza de la estadística con una estrella ·
calibración de imágenes con darks y flats · nubes sobre el sitio con el satélite GOES
\end{minipage}\par\vfill
{\large Diego Fernando Cagua Lozada\\ Fredy Esteban Cruz Ospina\\ Saray Juliana Pinto Sanabria}\par\vspace{0.4cm}
{\small\color{gris} Estudiantes del pregrado de Física · Departamento de Física\\
Universidad Nacional de Colombia, Bogotá D.C.\\[0.4em]
Desarrollado para la asignatura Técnicas Observacionales de la Maestría en Astronomía\\
Observatorio Astronómico Nacional\\[0.4em]
Interfaz y funcionalidades desarrolladas con ayuda de Claude (Anthropic)}\par\vspace{1cm}
{\small Noviembre de 2026 · \url{https://github.com/DiegoCagua/TurbuLENS}}
\end{titlepage}

\tableofcontents
\newpage

% ===================================================================================== 1
\section{Qué es TurbuLENS}

TurbuLENS es una aplicación de escritorio para Windows, macOS y Linux. Separa en imágenes los videos de un
\textbf{DIMM} (\emph{Differential Image Motion Monitor}, monitor de movimiento diferencial de imagen), mide en cada
cuadro la imagen doble de una estrella y calcula el \textbf{seeing}: qué tanto borronea la atmósfera las estrellas
en el sitio de observación. Además trae herramientas para enseñar estadística, calibrar imágenes astronómicas y
revisar si había nubes con el satélite GOES.

\begin{table}[H]
\centering\small
\begin{tabularx}{\textwidth}{@{}l Y@{}}
\toprule
Modo & Para qué sirve\\
\midrule
\textbf{Análisis científico} & Medir el seeing con un video del DIMM: recorte, tramo, parámetros del equipo,
procesamiento, resultados, informe PDF, historial y exportación de datos.\\
\textbf{Educativo} & Aprender estadística (histogramas, media, desviación estándar, la campana de Gauss, errores,
datos atípicos y correlación) con datos reales de una estrella, más tres pantallas de astronomía sencilla.\\
\textbf{Imágenes} & Calibrar imágenes con bias, darks y flats (archivos sueltos o carpetas completas, emparejados
según sus cabeceras) y leer el valor de cada píxel en un visor al estilo de SAOImage DS9.\\
\textbf{Satélite} & Descargar las imágenes del satélite GOES-Este para una fecha, hora y lugar, y ver canal por
canal si había nubes.\\
\bottomrule
\end{tabularx}
\caption{Los cuatro modos de TurbuLENS.}
\end{table}

\subsection{Descarga y apertura}

TurbuLENS no se instala. En la página de versiones del repositorio
(\url{https://github.com/DiegoCagua/TurbuLENS/releases/latest}) se descarga el archivo del sistema:

\begin{table}[H]
\centering\small
\begin{tabularx}{\textwidth}{@{}l l Y@{}}
\toprule
Sistema & Archivo & Cómo se abre\\
\midrule
Windows 10/11 & \codigo{TurbuLENS-v1.4.1-Windows-x64.zip} & Descomprimir y doble clic en \codigo{TurbuLENS.exe}.\\
Mac con Apple Silicon & \codigo{TurbuLENS-v1.4.1-macOS-AppleSilicon.zip} & Arrastrar \codigo{TurbuLENS.app} a Aplicaciones.\\
Mac con Intel & \codigo{TurbuLENS-v1.4.1-macOS-Intel.zip} & Igual que en Apple Silicon.\\
Linux & \codigo{TurbuLENS-v1.4.1-Linux-x64.tar.gz} & \codigo{chmod +x TurbuLENS} y doble clic.\\
\bottomrule
\end{tabularx}
\caption{Archivos de la versión 1.4.1.}
\end{table}

\begin{nota}{Primera apertura.}
TurbuLENS no tiene firma digital. En Windows puede aparecer «Windows protegió su PC»: se pulsa \emph{Más
información} y luego \emph{Ejecutar de todas formas}. En Mac, la primera vez se abre con clic derecho →
\emph{Abrir}. La primera apertura tarda unos segundos.
\end{nota}

\subsection{Idioma, tema y tamaño de la interfaz}

\begin{itemize}
\item \textbf{Idioma:} español, inglés o portugués. Se cambia en la pantalla de inicio o en \emph{Acerca de}; los
datos y las mediciones se conservan.
\item \textbf{Tema:} oscuro (naranja, para la noche de observación) o claro (verde azulado, para el aula y la
impresión). Se cambia con el sol o la luna de la barra superior.
\item \textbf{Tamaño de la interfaz} (nuevo en la 1.4.0): por defecto TurbuLENS se dibuja al \textbf{100\,\%}, aunque
Windows o Linux tengan la pantalla al 125\,\% o 150\,\%. Así las páginas caben completas en portátiles de
1366~×~768. En \emph{Acerca de} se puede elegir 125\,\%, 150\,\% o «Como el sistema»; el cambio se aplica al volver a
abrir la aplicación. En macOS siempre se respeta la pantalla Retina.
\end{itemize}

\captura{inicio.png}{Pantalla de inicio con los cuatro modos.}

% ===================================================================================== 2
\section{Fundamentos}

\subsection{La turbulencia atmosférica y el seeing}

La luz de una estrella llega a la Tierra como un frente de onda plano. Al cruzar la atmósfera atraviesa celdas de
aire de distinta temperatura, y por lo tanto de distinto índice de refracción, que se mueven con el viento. Cada
celda inclina un poco el frente de onda: la imagen de la estrella \textbf{tiembla}, cambia de brillo (titila) y, en
una exposición larga, se ensancha.

El tamaño de esas celdas se describe con el \textbf{parámetro de Fried}, $r_0$: el diámetro de una zona del frente de
onda en la que la turbulencia casi no lo deforma. Un cielo bueno tiene $r_0$ grande (10~cm o más en el visible). El
\textbf{seeing} $\varepsilon$ es el ancho a media altura (FWHM) de la imagen de una estrella en exposición larga:
\begin{equation}
\varepsilon = 0{,}98\,\frac{\lambda}{r_0}\qquad\text{(en radianes; $\times 206\,265$ para segundos de arco)}.
\label{eq:seeing}
\end{equation}

\begin{table}[H]
\centering\small
\begin{tabular}{@{}l l@{}}
\toprule
Seeing & Calidad del cielo\\
\midrule
$\lesssim 1''$ & Excelente (los mejores observatorios)\\
$2''$ a $3''$ & Aceptable\\
$> 3''$ & Muy turbulento\\
\bottomrule
\end{tabular}
\caption{Escala de calidad que usa TurbuLENS. $1''$ es el tamaño de una moneda de 1~cm vista a 2~km.}
\end{table}

\subsection{El DIMM}

Medir el seeing con el tamaño de la imagen de una estrella no funciona bien con telescopios pequeños: el foco, las
vibraciones y el seguimiento también la ensanchan. El DIMM evita esos problemas. Sobre la apertura del telescopio se
pone una \textbf{máscara con dos orificios} de diámetro $D$ separados una distancia $B$ entre centros. La luz entra
por dos caminos y la cámara, ligeramente desenfocada, ve \textbf{dos imágenes} de la misma estrella.

\begin{itemize}
\item Si el telescopio tiembla o el seguimiento falla, las dos imágenes se mueven \textbf{juntas}: la distancia entre
ellas no cambia.
\item La turbulencia, en cambio, inclina distinto el frente de onda en cada orificio: las imágenes se acercan y se
alejan \textbf{entre sí}.
\end{itemize}

Por eso el DIMM no mide dónde está la estrella sino la \textbf{distancia entre las dos imágenes}, cuadro a cuadro. La
varianza de esa distancia depende solo de la turbulencia.

\captura{edu_dimm.png}{La pantalla «Los dos huequitos» del modo educativo: con el temblor los puntos se mueven juntos;
con la turbulencia cambia la distancia.}

\subsection{Del movimiento diferencial al seeing: el modelo de Tokovinin}

Sea $\vec d$ el vector que une los dos centroides en un cuadro. Su movimiento se separa en dos componentes: la
\textbf{longitudinal} $\ell$, a lo largo de la línea que une los orificios, y la \textbf{transversal} $t$,
perpendicular. Tokovinin (2002) dedujo, para el centroide de cada imagen,
\begin{align}
\sigma_\ell^2 &= K_\ell\,\lambda^2\,r_0^{-5/3}\,D^{-1/3}, & K_\ell &= 0{,}364\left(1 - 0{,}532\,b^{-1/3} - 0{,}024\,b^{-7/3}\right),\\
\sigma_t^2 &= K_t\,\lambda^2\,r_0^{-5/3}\,D^{-1/3}, & K_t &= 0{,}364\left(1 - 0{,}798\,b^{-1/3} + 0{,}018\,b^{-7/3}\right),
\end{align}
con $b = B/D$ y las varianzas en radianes cuadrados. Despejando,
\begin{equation}
r_0 = \left(\frac{K\,\lambda^2\,D^{-1/3}}{\sigma^2}\right)^{3/5},
\end{equation}
y con la ecuación~\eqref{eq:seeing} sale un seeing longitudinal $\varepsilon_\ell$ y uno transversal $\varepsilon_t$.
TurbuLENS reporta su promedio. Las varianzas medidas en píxeles se pasan a radianes con la escala de placa $s$:
\begin{equation}
\sigma^2\,[\text{rad}^2] = \sigma^2\,[\text{px}^2]\left(\frac{s}{206\,264{,}8}\right)^2 .
\end{equation}

\subsubsection{Corrección al cenit}

Cerca del horizonte la luz cruza más aire. Si se conoce la altura $h$ de la estrella, con $z = 90^\circ - h$,
\begin{equation}
\varepsilon_{\text{cenit}} = \varepsilon\,(\cos z)^{3/5},
\end{equation}
que permite comparar noches y estrellas distintas. La altura se obtiene, por ejemplo, en Stellarium con la fecha, la
hora y el lugar de la grabación.

\subsection{La escala de placa}\label{sec:escala}

La escala de placa dice cuántos segundos de arco del cielo caben en un píxel:
\begin{equation}
s\,[''/\text{px}] = \frac{206{,}265 \times p\,[\text{µm}]}{f\,[\text{mm}]},
\end{equation}
donde $p$ es el tamaño del píxel y $f$ la focal del telescopio. Es el dato que más pesa en el resultado: como
$\sigma^2 \propto s^2$ y $\varepsilon \propto r_0^{-1} \propto (\sigma^2)^{3/5}$,
\begin{equation}
\varepsilon \propto s^{6/5}.
\end{equation}
Un error de un factor 3 en la escala da un seeing casi 4 veces menor o mayor.

\begin{cuidado}{El píxel efectivo del video.}
Las cámaras de fotos (Nikon, Canon) y los celulares \textbf{reducen} la imagen completa del sensor para grabar
video. Una Nikon D5600 tiene 6000 píxeles de ancho, pero su video de 1920 px usa todo el ancho del sensor: cada
píxel del video abarca unos 3 píxeles del sensor. En video vale el \textbf{píxel efectivo}:
\[
p_{\text{efectivo}} = \frac{\text{ancho del sensor}}{\text{ancho del video}} = \frac{23{,}5\ \text{mm}}{1920\ \text{px}} = 12{,}24\ \text{µm}
\quad(\text{no } 3{,}9\ \text{µm}).
\]
Con el Celestron NexStar 8SE ($f = 2032$~mm) la escala correcta es $1{,}242''/\text{px}$. Con el píxel del sensor
saldría $0{,}398''/\text{px}$ y el seeing, unas 3,9 veces menor (por ejemplo $\approx 1''$ en vez de
$\approx 3{,}8''$). Las cámaras astronómicas (QHY, ZWO) graban con los píxeles reales del sensor.
\end{cuidado}

Lo más exacto es medir la escala con una estrella doble de separación conocida, grabada con el mismo equipo.

\subsection{Cómo se miden los centroides}\label{sec:centroides}

Cada cuadro del tramo elegido se procesa así (en bloques vectorizados, sin recorrer los píxeles uno a uno):
\begin{enumerate}
\item \textbf{Fondo y ruido:} el fondo es la mediana del recorte y el ruido, $1{,}4826\times$ la desviación absoluta
mediana (MAD). Si hay maestros de calibración, el cuadro se corrige antes como (cuadro $-$ dark) $\div$ flat.
\item \textbf{Píxeles con luz:} se suaviza la imagen con una caja de $3\times3$ y se toman los píxeles que superan
$\max(0{,}15\times\text{pico},\ 4\times\text{ruido}/3)$.
\item \textbf{Dos grupos:} esos píxeles se separan en dos grupos con k-medias ponderado por la intensidad
(8 iteraciones), empezando en el píxel más brillante y en el píxel con luz más alejado de él. Así se separan bien
spots pequeños, grandes, en anillo (desenfocados) o partidos.
\item \textbf{Centroides:} el centro de cada spot es el centroide ponderado por intensidad de su grupo, recortado a
$3r + 3$ píxeles (con $r$ el radio de giro). Su precisión es de centésimas de píxel.
\item \textbf{Validación:} el cuadro cuenta si los dos spots tienen suficiente señal sobre el ruido, el segundo tiene
al menos el 12\,\% de la luz del primero, están separados más de $2{,}2$ veces su radio, hay una caída de luz entre
ellos (menos del 60\,\% del pico menor en el punto medio) y ninguno toca el borde del recorte.
\item \textbf{Atípicos:} se descartan los cuadros cuya separación está a más de 5 desviaciones robustas de la mediana.
\end{enumerate}

\textbf{Seguimiento:} antes de medir, TurbuLENS ubica la estrella aunque la región haya quedado lejos (la búsqueda se
amplía hasta el cuadro completo) y, con \emph{Tamaño automático}, ajusta la región a la separación y al tamaño de los
spots. Con \emph{Seguir a la estrella} la región se recentra en cada cuadro; si la estrella se pierde 5 cuadros
seguidos, se vuelve a buscar en todo el cuadro.

\subsection{Controles de calidad: las correlaciones}\label{sec:correlacion}

El coeficiente de correlación de Pearson
\begin{equation}
r = \frac{\sum_i (x_i - \bar x)(y_i - \bar y)}{\sqrt{\sum_i (x_i - \bar x)^2\,\sum_i (y_i - \bar y)^2}}
\end{equation}
va de $-1$ a $1$. TurbuLENS calcula tres controles, que aparecen en Resultados y en el informe PDF:
\begin{itemize}
\item \textbf{$r$ longitudinal–transversal:} la turbulencia mueve las dos componentes de forma independiente, así que
debe estar cerca de 0 ($|r| < 0{,}2$). Un $|r|$ grande delata vibración o viento en una sola dirección, cambios de
foco, una línea de base mal estimada o spots confundidos.
\item \textbf{$r$ entre los dos spots} (longitudinal y transversal, tras quitar la deriva lenta con una media móvil de
1~s): cerca de 1 domina el movimiento común (vibración, seguimiento), que el DIMM cancela; entre 0 y 1 es lo normal;
cerca de 0, cada orificio ve turbulencia independiente; negativo indica cambios de foco o de escala, o errores de
detección.
\item \textbf{$\sigma_t/\sigma_\ell$} medido frente al teórico $\sqrt{K_t/K_\ell}$: si se parecen, la turbulencia se
comporta como predice el modelo.
\end{itemize}

\subsection{Calibración de imágenes: bias, darks y flats}\label{sec:calibracion}

Una imagen cruda de una cámara astronómica o de una DSLR mezcla la luz del cielo con defectos del instrumento:

\begin{table}[H]
\centering\small
\begin{tabularx}{\textwidth}{@{}l Y Y@{}}
\toprule
Imagen & Cómo se toma & Qué corrige\\
\midrule
\textbf{Bias} & La exposición más corta posible, con la tapa puesta. & El nivel de lectura de la cámara (offset).\\
\textbf{Dark} & Con la tapa puesta, con la \textbf{misma exposición, ganancia (o ISO) y temperatura} que las luces. &
El bias más la corriente oscura y los píxeles calientes. Se resta.\\
\textbf{Flat} & De una superficie iluminada de forma pareja (cielo del atardecer, pantalla blanca), con el
\textbf{mismo foco y filtro}. & El viñeteo y las sombras del polvo. Se divide.\\
\textbf{Dark de flat} & Darks con la exposición de los flats. & Lo que se resta a cada flat antes de combinarlos.\\
\bottomrule
\end{tabularx}
\caption{Imágenes de calibración.}
\end{table}

\textbf{Por qué hay que emparejarlas.} La corriente oscura crece con el tiempo de exposición (aproximadamente
$\text{dark}(t) \approx \text{bias} + I_d\,t$) y con la temperatura del sensor (se duplica cada 5–7~°C en muchos
sensores), y su escala cambia con la ganancia. Un dark de 30~s no sirve para una luz de 120~s. El flat, en cambio,
depende del filtro, del foco y de la posición de la cámara. Por eso TurbuLENS empareja cada luz con los darks de su
misma exposición, ganancia, temperatura, binning y tamaño, y con los flats de su mismo filtro.

\textbf{Maestros.} Cada grupo se combina en una imagen maestra: la \textbf{mediana} (predeterminada) es robusta a los
rayos cósmicos; la \textbf{media con recorte de 3\,σ} descarta en cada píxel los valores a más de 3 desviaciones
robustas; la \textbf{media} reduce más el ruido si no hay valores raros. Cada flat se nivela a mediana 1 antes de
combinarse, y el flat maestro se normaliza a mediana 1 (los píxeles muertos quedan en 1).

\textbf{Corrección:}
\begin{equation}
\text{calibrada} = \frac{\text{luz} - \text{dark}}{\text{flat normalizado}} .
\end{equation}

\textbf{Dark escalado.} Si no hay darks de la exposición de una luz pero sí bias, TurbuLENS escala un dark de otra
exposición (prefiere uno más largo, para no amplificar su ruido):
\begin{equation}
\text{dark}(t_{\text{luz}}) = \left(\text{dark}(t_{\text{dark}}) - \text{bias}\right)\frac{t_{\text{luz}}}{t_{\text{dark}}} + \text{bias}.
\end{equation}
Si no hay bias, se usa el dark sin escalar o solo el bias, y el plan lo advierte.

\subsection{El satélite GOES}\label{sec:goes}

Los satélites geoestacionarios GOES de la NOAA miran América desde 35\,786~km de altura. Su instrumento ABI toma una
imagen de \textbf{disco completo} cada 10 minutos en 16 canales. TurbuLENS usa el producto
\codigo{ABI-L2-CMIPF} (imágenes ya calibradas): en los canales 1 a 6 el valor es la reflectancia (útil solo de
día); en los canales 7 a 16, la temperatura de brillo en kelvin, de día y de noche. Para nubes de noche el más usado
es el \textbf{canal 13} (10,3~µm): entre más fría, más alta la nube.

\begin{table}[H]
\centering\small
\begin{tabular}{@{}l l@{}}
\toprule
Temperatura en el canal 13 & Lectura\\
\midrule
$< -40$~°C & Topes de nubes altas y frías\\
$-40$ a $-10$~°C & Nubes medias o altas\\
$-10$ a $5$~°C & Posibles nubes bajas o neblina\\
$> 5$~°C & Probablemente despejado (se ve la superficie)\\
\bottomrule
\end{tabular}
\caption{Lectura sencilla de los canales infrarrojos limpios (13, 14 y 15).}
\end{table}

\textbf{Qué satélite.} GOES-16 fue el GOES-Este hasta el 4 de abril de 2025 a las 15:10~UTC; desde entonces lo es
GOES-19. TurbuLENS elige solo según la fecha y, si el archivo no está en un repositorio, lo busca en el otro.

\textbf{Del lugar al píxel.} La latitud y la longitud se convierten en los ángulos de barrido del instrumento con
la proyección geoestacionaria de la guía de usuario de GOES-R (elipsoide GRS80, satélite en $-75^\circ$ de longitud
para GOES-Este), y de ahí en la fila y la columna de la imagen. El píxel es de unos 2~km en los canales infrarrojos.

% ===================================================================================== 3
\section{Equipos del OAN}\label{sec:oan}

El botón \boton{Dispositivos del OAN…} de Parámetros (sección~\ref{sec:parametros}) llena $D$, $B$ y la escala de
placa con estos equipos:

\begin{table}[H]
\centering\small
\begin{tabularx}{\textwidth}{@{}l r r Y@{}}
\toprule
Telescopio & Apertura & Focal & Notas\\
\midrule
Celestron NexStar 8SE & 203~mm & 2032~mm (f/10) & Schmidt-Cassegrain; el de las mediciones del grupo.\\
Meade 10″ (LX200) & 254~mm & 2500~mm (f/10) & Comprobar la focal en la etiqueta: con reductor f/6,3 es de
unos 1600~mm.\\
\bottomrule
\end{tabularx}
\caption{Telescopios.}
\end{table}

\begin{table}[H]
\centering\small
\begin{tabularx}{\textwidth}{@{}l l r Y@{}}
\toprule
Cámara & Sensor & Píxel & En video\\
\midrule
Nikon D5600 & APS-C 23,5~×~15,6~mm, 6000~×~4000 & 3,92~µm & Reduce el sensor a 1920~px: píxel efectivo 12,24~µm.\\
Canon T5i (700D) & APS-C 22,3~×~14,9~mm, 5184~×~3456 & 4,30~µm & Reduce el sensor a 1920~px: píxel efectivo 11,61~µm.\\
QHY174M-GPS & IMX174 11,25~×~7,03~mm, 1920~×~1200 & 5,86~µm & Graba con los píxeles reales.\\
\bottomrule
\end{tabularx}
\caption{Cámaras.}
\end{table}

\begin{table}[H]
\centering\small
\begin{tabular}{@{}l c c c@{}}
\toprule
 & Nikon D5600 (video 1920) & Canon T5i (video 1920) & QHY174M-GPS\\
\midrule
Celestron NexStar 8SE & 1,242 & 1,179 & 0,595\\
Meade 10″ (2500~mm) & 1,010 & 0,958 & 0,483\\
\bottomrule
\end{tabular}
\caption{Escalas de placa resultantes, en $''/\text{px}$.}
\end{table}

\textbf{Máscaras:} la del grupo para el 8SE tiene orificios de $D = 4{,}2$~cm separados $B = 14{,}4$~cm; la del
protocolo de ANDEAN ROAD para 8 pulgadas, $D = 2{,}5$~cm y $B = 14$~cm. Para otra máscara se elige
\emph{Conservar D y B} y se escriben a mano.

\captura{oan.png}{Ventana «Dispositivos del OAN».}

% ===================================================================================== 4
\section{Uso: modo análisis}

El menú de la izquierda muestra los pasos numerados. Se puede ir y volver entre ellos sin perder nada.

\subsection{1. Video}

Se arrastra el video a la ventana o se busca con el selector de archivos del sistema. TurbuLENS abre cientos de
formatos (MOV, MP4, AVI, MKV, WMV, MTS, WebM, SER, FITS…) con FFmpeg incluido, también los de 16 bits de cámaras
astronómicas. El video no se copia: se lee desde su carpeta. A la derecha aparece su información.

\captura{video.png}{Video cargado.}

\subsection{2. Recorte y tramo}

\textbf{Región.} Se dibuja como en Paint alrededor de la estrella (arrastrar, mover, cambiar el tamaño con las
manijas, rueda para el zoom) o con \boton{◎ Ajustar a la estrella}. A la derecha se ve el interior de la región y un
estado: «Contiene los dos puntos ✓», muy ajustada o muy grande. Las flechas del teclado mueven la región un píxel y
Mayús~+~flechas cambian su tamaño.

\begin{itemize}
\item \textbf{Seguir a la estrella:} recentra la región en cada cuadro (necesario si el telescopio no tiene motor de
seguimiento). Al reproducir, la región sigue a la estrella en vivo.
\item \textbf{Tamaño automático:} ajusta el tamaño de la región a los spots al procesar. Si la región quedó muy
pequeña, se agranda sola.
\end{itemize}

\textbf{Tramo.} En la línea de tiempo se marcan la entrada y la salida con \boton{[ Entrada} y \boton{Salida ]}, o se
arrastran sus marcas. Los botones \codigo{<<} y \codigo{>>} reproducen a velocidad constante (0,25× a 32× respecto al
video real), \codigo{<} y \codigo{>} avanzan un cuadro, y \emph{Ir a} salta a un cuadro o a un tiempo
\codigo{mm:ss.s}. Se recomiendan al menos 15 segundos.

\captura{region.png}{Recorte y tramo.}

\subsection{3. Parámetros}\label{sec:parametros}

\begin{itemize}
\item \boton{Dispositivos del OAN…} llena todo con los equipos del observatorio (sección~\ref{sec:oan}).
\item \textbf{D}, \textbf{B} y \textbf{λ} con unidades a elegir (µm, mm, cm o m; la unidad se recuerda). Los números
aceptan punto o coma decimal.
\item \textbf{Escala de placa.} \boton{Calcular…} abre la calculadora: cámara (lista de cámaras frecuentes), casilla
\emph{Grabé en video} (usa el píxel efectivo, con el ancho tomado del video cargado), focal y \boton{Usar}.
\item \textbf{Altura de la estrella} (opcional): para el seeing en el cenit.
\item \textbf{Perfil:} \boton{Guardar…} crea un perfil con los valores actuales; el botón \boton{⋯} los actualiza,
renombra o elimina. Los perfiles son los mismos en los modos análisis y educativo.
\item \textbf{Imágenes a guardar:} PNG (recomendado), TIFF de 16 bits, JPG, cubo NumPy o cubo FITS; 1 de cada $n$
cuadros; carpeta de salida y estimación del espacio en disco.
\item \textbf{Calibración con dark y flat:} al activarla aparecen los campos del dark, flat y bias maestros (se crean
en el modo Imágenes).
\end{itemize}

Si ya hay resultados, los cambios de $D$, $B$, λ, escala y altura se aplican al instante: no hace falta procesar de
nuevo.

\captura{parametros.png}{Parámetros después de usar «Dispositivos del OAN».}

\subsection{4. Procesar}

Un resumen dice qué se va a hacer. \boton{▶ Iniciar} empieza; la barra muestra cuadros por segundo y el tiempo que
falta, y \boton{Cancelar} detiene sin perder lo medido. \textbf{Sin escala de placa no se puede procesar}: un aviso
lo explica y lleva a Parámetros.

\captura{procesar_sin_escala.png}{Aviso cuando falta la escala de placa.}

\subsection{5. Resultados}

Arriba van los indicadores: seeing estimado (con su calidad en palabras y color), longitudinal y transversal, $r_0$,
separación media, desviación estándar, cuadros usados y seeing en el cenit. Cada uno tiene su botón \boton{!}. Debajo:
\begin{itemize}
\item la \textbf{animación de los centroides}, que muestra cuadro a cuadro dónde ubicó TurbuLENS cada punto;
\item el \textbf{histograma} de la separación con la curva normal y la media;
\item la \textbf{σ en el tiempo} (una subida aislada suele ser un golpe o una ráfaga de viento);
\item el \textbf{diagrama del movimiento diferencial} con las correlaciones $r$ (sección~\ref{sec:correlacion}).
\end{itemize}
Con \boton{Guardar en historial} (con una etiqueta opcional) el análisis pasa al historial.

\captura{resultados.png}{Resultados.}

\textbf{Exportar.} Solo se elige la carpeta; el nombre de los archivos es automático y nunca se sobrescribe nada:
\begin{itemize}
\item datos en CSV, Excel, JSON, NumPy y tabla FITS;
\item gráficas en PDF, PNG y SVG;
\item \textbf{informe PDF} con el resumen, las gráficas y la sección «Cómo leer la correlación $r$» con los valores
del análisis y un diagnóstico.
\end{itemize}

\subsection{6. Historial}

Una base de datos local guarda cada análisis para comparar noches y épocas del año, con etiquetas y notas. Es el
mismo en los modos análisis y educativo (una columna dice en cuál se hizo). Se exporta igual que los resultados.

\captura{historial.png}{Historial.}

% ===================================================================================== 5
\section{Uso: modo educativo}

El modo educativo usa la estrella como pretexto para aprender estadística: cada cuadro del video es un dato.
\begin{itemize}
\item \textbf{La estrella} (va primero): «¿Por qué titilan?» (con la animación de la moneda en la piscina y la luz
que cruza el aire), «Los dos huequitos» (la máscara y el telescopio que tiembla) y un simulador.
\item \textbf{Estadística}: nueve lecciones (datos y mediciones, histograma, media y mediana, desviación estándar,
campana de Gauss, muestras y error, datos atípicos, correlación y laboratorio libre) y una prueba final de 11
preguntas.
\item \textbf{Tu video}: el mismo flujo del modo análisis, con explicaciones adicionales.
\end{itemize}
Cada lección tiene una explicación corta, la fórmula, «Pruébalo», una gráfica que responde a los controles, una tabla
de los datos mostrados (se guarda en CSV o Excel) y una pregunta que se responde una sola vez por sesión, con las
opciones en orden al azar. Los datos pueden ser los de ejemplo (Altair) o los del video que procese la clase.

\captura{edu_titilan.png}{«¿Por qué titilan las estrellas?» con su animación.}
\captura{edu_l3.png}{Lección de media y mediana: los datos raros caen a un lado elegido al azar.}

% ===================================================================================== 6
\section{Uso: modo Imágenes}

\subsection{Visor de píxeles}

Abre imágenes FITS (también cubos), PNG, TIFF de 16 bits, JPG y NumPy, y funciona como SAOImage DS9:
\begin{itemize}
\item \textbf{Rueda:} acerca o aleja alrededor del cursor. \textbf{Arrastrar:} mueve la imagen.
\item Bajo la imagen se lee la \textbf{posición} y el \textbf{valor} del píxel bajo el cursor (columna y fila desde 1;
la fila 1 es la de arriba).
\item \textbf{Escala:} lineal, logarítmica, raíz cuadrada, asinh o histograma ecualizado. \textbf{Límites:} mínimo–máximo,
percentiles o ZScale (el de DS9). \textbf{Colores:} grises, invertidos, calor, viridis o frío–cálido. Los datos no
cambian: solo cómo se ven.
\item \textbf{Región:} Mayús~+~arrastrar (o clic derecho y arrastrar) calcula píxeles, media, mediana, σ, σ robusta,
mínimo, máximo y suma. \boton{Exportar CSV…} guarda x, y y el valor de cada píxel de la región.
\item La \textbf{cabecera FITS} se ve a la derecha; la imagen se guarda en FITS de 32 bits o en PNG de 16 bits.
\end{itemize}

\captura{visor.png}{Visor de píxeles con una región.}

\subsection{Calibración con darks y flats}

Funciona con \textbf{imágenes sueltas} o con \textbf{carpetas completas} (se recorren también las subcarpetas):
\begin{enumerate}
\item \boton{Agregar carpeta…} o \boton{Agregar imágenes…}. TurbuLENS lee la \textbf{cabecera} de cada archivo sin
cargar la imagen: tipo (\codigo{IMAGETYP}, \codigo{FRAME}…), exposición (\codigo{EXPTIME}, \codigo{EXPOSURE}), ganancia
o ISO, temperatura (\codigo{CCD-TEMP}), filtro, binning, tamaño y cámara. En TIFF y JPG de cámaras de fotos lee los
datos EXIF (exposición e ISO).
\item Cada archivo se \textbf{clasifica} en luz, dark, flat, dark de flat o bias: por la cabecera; si no la tiene, por el
nombre del archivo («dark\_30s\_001») o de la carpeta («Darks/», «Flats/»). La columna \emph{Según} dice de dónde salió;
lo que no se pudo saber se toma como luz y se marca en amarillo. Si alguna quedó mal, se selecciona y se cambia con
\emph{márcala como}.
\item \textbf{Qué se va a hacer:} un plan agrupa las luces por exposición, filtro, ganancia y tamaño, y dice para cada
grupo qué se resta (dark de la misma exposición, dark escalado con bias, solo bias o nada) y qué se divide (flat del
mismo filtro), con avisos en amarillo.
\item \boton{Crear maestros y calibrar}: cada maestro se crea una sola vez y se guarda en \codigo{maestros/}; cada luz
calibrada, en \codigo{calibradas/}, con su cabecera original más \codigo{CALSTAT}, \codigo{NDARK}, \codigo{NFLAT} y, si se
escaló el dark, \codigo{DARKSCAL}. Las series grandes se combinan por franjas con un archivo temporal, sin llenar la
memoria. Se puede cancelar.
\item \boton{Usar en el análisis de video} pasa los maestros del grupo principal a la calibración del modo análisis.
\end{enumerate}

\captura{calibracion_plan.png}{Una noche con luces de 30 y 60~s en dos filtros: el plan escala el dark de 120~s para
las luces de 60~s y usa el flat de cada filtro.}

\begin{nota}{Si solo hay calibraciones.}
Sin luces, TurbuLENS crea solo los maestros (bias, darks por exposición y flats por filtro), listos para otra noche.
\end{nota}

% ===================================================================================== 7
\section{Uso: modo Satélite}

\begin{enumerate}
\item \textbf{Fecha, hora y lugar:} la hora puede ser de Colombia (UTC$-5$) o UTC; las coordenadas por defecto son
las del OAN (4,6386° N, 74,0846° O). \emph{Satélite}: automático según la fecha, GOES-16 o GOES-19.
\item \textbf{Canales:} del 1 al 16 (al pasar el ratón se ve para qué sirve cada uno y cuánto pesa su archivo).
\emph{Solo el 13} basta para saber si había nubes.
\item \textbf{Carpeta de descarga} y \boton{Buscar, descargar y leer}: se descarga la imagen de disco completo más
cercana a la hora pedida (≈25~MB por canal infrarrojo; el canal 2 pesa ≈330~MB). Si ya está en la carpeta, no se
descarga otra vez. \emph{Leer archivos ya descargados} abre archivos \codigo{.nc} que ya se tengan.
\end{enumerate}

La tabla muestra para cada canal el valor en el punto, el promedio de $3\times3$ píxeles, una lectura sencilla y la
hora real del escaneo. Debajo se ve un recorte alrededor del punto (en infrarrojo las nubes frías se ven blancas),
que se puede abrir en el visor de píxeles. En la carpeta quedan un CSV con los valores y los recortes en FITS y PNG.

\captura{satelite.png}{Modo satélite con una lectura del canal 13.}

\begin{nota}{Datos.}
NOAA GOES-R Series, producto ABI L2 Cloud and Moisture Imagery (CMIPF), del repositorio abierto NOAA Open Data
Dissemination en Amazon Web Services. Hace falta conexión a internet para descargar.
\end{nota}

% ===================================================================================== 8
\section{Archivos que produce TurbuLENS}

\begin{table}[H]
\centering\small
\begin{tabularx}{\textwidth}{@{}l Y@{}}
\toprule
Archivo & Contenido\\
\midrule
\codigo{cuadro\_000000.png} … & Cada cuadro del tramo, recortado a la región (o TIFF, JPG, un cubo \codigo{.npy} o \codigo{.fits}).\\
\codigo{mediciones.csv} & Por cuadro: tiempo, posiciones de los dos spots, flujos, separación y si es válido.\\
\codigo{resumen.json} & Parámetros, opciones del procesamiento, seeing y estadísticas.\\
Exportación de resultados & CSV, Excel (resumen y mediciones), JSON, NumPy (\codigo{.npz}) y tabla FITS; gráficas en PDF, PNG y SVG; informe PDF.\\
Historial & Base SQLite local en la carpeta de datos del usuario; se exporta a CSV o Excel con gráficas.\\
\codigo{maestros/maestro\_*.fits} & Bias, darks y flats maestros (FITS de 32 bits, con \codigo{IMAGETYP}, \codigo{NCOMBINE}, \codigo{EXPTIME}, \codigo{FILTER}).\\
\codigo{calibradas/*\_calibrada.fits} & Luces calibradas con su cabecera original más los datos de la calibración.\\
\codigo{turbulens\_goes\_*\_valores.csv} & Valores de cada canal en el punto.\\
\codigo{turbulens\_goes\_*\_C13\_recorte.fits} & Recorte de cada canal (también en PNG).\\
\bottomrule
\end{tabularx}
\caption{Archivos de salida.}
\end{table}

Cada archivo de datos, gráfica e informe lleva una línea de crédito con la cita de TurbuLENS.

% ===================================================================================== 9
\section{Solución de problemas}

\begin{table}[H]
\centering\small
\begin{tabularx}{\textwidth}{@{}Y Y@{}}
\toprule
Problema & Qué hacer\\
\midrule
El seeing sale imposiblemente bajo (≈1″) con una Nikon o Canon. & La escala se calculó con el píxel del sensor. Marca
\emph{Grabé en video} en la calculadora o usa \boton{Dispositivos del OAN…} (sección~\ref{sec:escala}).\\
No deja procesar. & Falta la escala de placa: escríbela en Parámetros.\\
«No se ven los dos puntos» o pocos cuadros válidos. & Agranda la región o activa \emph{Tamaño automático}; revisa que
el video tenga las dos imágenes separadas (desenfoca un poco menos si se tocan).\\
La región se pierde con el tiempo. & Activa \emph{Seguir a la estrella}.\\
Una subida aislada de σ en el tiempo. & Un golpe o una ráfaga: ajusta el tramo para dejarla fuera.\\
$|r|$ longitudinal–transversal mayor que 0,2. & Hay un efecto que no es turbulencia (vibración en una dirección,
foco): revisa el montaje y el tramo.\\
La interfaz se ve muy grande o no cabe. & En \emph{Acerca de} elige el tamaño 100\,\% y vuelve a abrir TurbuLENS.\\
El plan de calibración dice «sin darks» para un grupo. & No hay darks de su tamaño, binning o ganancia; o falta el
bias para escalar uno de otra exposición.\\
Una imagen quedó como «supuesto». & No tiene \codigo{IMAGETYP} ni palabras clave en el nombre: márcala a mano.\\
«No coincide con el video» en la calibración del análisis. & Los maestros deben medir lo mismo que el video
(ancho y alto).\\
El modo satélite no encuentra imágenes. & Revisa la conexión a internet, la fecha (GOES-16 desde diciembre de 2017) y
la hora.\\
Windows muestra «Windows protegió su PC». & \emph{Más información} → \emph{Ejecutar de todas formas}.\\
\bottomrule
\end{tabularx}
\caption{Problemas frecuentes.}
\end{table}

% ===================================================================================== 10
\section{Historial de versiones}

\begin{table}[H]
\centering\small
\begin{tabularx}{\textwidth}{@{}l Y@{}}
\toprule
Versión & Cambios principales\\
\midrule
1.0.0 & Primera versión: extracción de cuadros, medición vectorizada, seeing con Tokovinin, modo educativo, historial.
Licencia MIT.\\
1.1.x & Animación de los centroides; arreglos en Linux; Fredy Esteban Cruz Ospina como coautor; licencia PolyForm
Noncommercial.\\
1.2.0–1.2.2 & Tres idiomas; los mismos controles de reproducción en toda la app; archivos para descomprimir y abrir;
historial y perfiles compartidos entre modos; exportar con nombres automáticos.\\
1.3.0 & Detector de dos grupos (spots grandes y en anillo); seguimiento cuadro a cuadro y tamaño automático; sin escala
no se procesa; calculadora con píxel efectivo de video; unidades; perfiles editables; informe con la explicación de
$r$; modos Imágenes y Satélite; animaciones del modo educativo; Saray Juliana Pinto Sanabria como coautora y tester.\\
\textbf{1.4.0} & Calibración guiada por las cabeceras para imágenes sueltas y carpetas completas (empareja darks por
exposición, ganancia, temperatura y binning, y flats por filtro; escala darks con el bias); botón \emph{Dispositivos
del OAN}; interfaz al 100\,\% por defecto; Parámetros y Recorte más limpios; este manual.\\
1.4.1 & Créditos: la interfaz y las funcionalidades se desarrollaron con ayuda de Claude (Anthropic).\\
\bottomrule
\end{tabularx}
\caption{Versiones de TurbuLENS.}
\end{table}

% ===================================================================================== 11
\section{Créditos, licencia y cómo citar}

\textbf{Autores:} Diego Fernando Cagua Lozada, Fredy Esteban Cruz Ospina y Saray Juliana Pinto Sanabria (coautora y
pruebas), estudiantes del pregrado de Física, Departamento de Física, Universidad Nacional de Colombia, Bogotá D.C.
Desarrollado para la asignatura Técnicas Observacionales de la Maestría en Astronomía del Observatorio Astronómico
Nacional.

\textbf{Desarrollo con inteligencia artificial:} La interfaz y las funcionalidades se desarrollaron con ayuda de Claude (Anthropic), un asistente de inteligencia artificial, bajo la dirección y revisión de los autores.

\textbf{Licencia:} PolyForm Noncommercial 1.0.0. Se puede usar, estudiar, modificar y compartir con fines no comerciales
(estudio, investigación, enseñanza), conservando la licencia y el crédito a los autores. No se permite el uso
comercial sin permiso escrito. Los componentes de terceros (Qt/PySide6, FFmpeg/PyAV, NumPy, Pillow, openpyxl, pyfive)
conservan sus licencias.

\textbf{Cómo citar:}
\begin{quote}\small
Cagua Lozada, D. F., Cruz Ospina, F. E., \& Pinto Sanabria, S. J. (2026). \emph{TurbuLENS} (versión 1.4.1)
[Software]. Observatorio Astronómico Nacional, Universidad Nacional de Colombia.
\url{https://github.com/DiegoCagua/TurbuLENS}
\end{quote}

\captura{acerca.png}{Acerca de: créditos, cita, idioma y tamaño de la interfaz.}

\section*{Referencias}
\addcontentsline{toc}{section}{Referencias}
\begin{enumerate}[label={[\arabic*]},leftmargin=*]
\item A. Tokovinin, «From differential image motion to seeing», \emph{Publications of the Astronomical Society of the
Pacific} \textbf{114}, 1156 (2002). \url{https://arxiv.org/abs/astro-ph/0111022}
\item M. Sarazin y F. Roddier, «The ESO differential image motion monitor», \emph{Astronomy \& Astrophysics}
\textbf{227}, 294 (1990).
\item D. L. Fried, «Optical resolution through a randomly inhomogeneous medium for very long and very short
exposures», \emph{Journal of the Optical Society of America} \textbf{56}, 1372 (1966).
\item ANDEAN ROAD, Oficina Regional Andina de la Astronomía para el Desarrollo, \emph{Determina la nitidez del cielo
nocturno en tu lugar de observación} (protocolo, 2026).
\item NOAA/NASA, \emph{GOES-R Series Product Definition and Users' Guide}, vol. 3 (Level 1b) y vol. 5 (Level 2+).
\item NOAA, \emph{GOES-19 operational transition to GOES-East}, aviso del 4 de abril de 2025.
\end{enumerate}

\end{document}
